\documentclass[tikz,border=12pt]{standalone}
\usetikzlibrary{decorations.pathmorphing}
\definecolor{raspberry}{rgb}{0.65,0.2,0.4}
\begin{document}
{ % begin diagram
    % ===== Parameters you can tweak =====
    % Cell size (increase these to make characters larger)
    \newcommand{\cellw}{0.38cm}
    \newcommand{\cellh}{0.38cm}

    % Cross-cut position (how much of left/top to keep)
    \newcommand{\cutX}{25}   % show 0..cutX columns on the left
    \newcommand{\cutY}{15}   % show 0..cutY rows on the top

    % Cutline params
 %   0.35mm, segment length=3mm
    \newcommand{\amp}{0.09}   % amplitude (cells)
    \newcommand{\lam}{1}    % wavelength (cells) 0.789
    \newcommand{\step}{0.25}  % sampling step (cells): smaller = smoother
    \newcommand{\samps}{400}  % smoothness

    % Corner patch sizes on the far sides
    \newcommand{\rightW}{3}  % width of right-side patches (near col 79)
    \newcommand{\bottomH}{3} % height of bottom-side patches (near row 59)

    % Tiny visual gaps (kept small even if many cells omitted)
    \newcommand{\gapX}{0.5}
    \newcommand{\gapY}{0.5}

    % Size of the compressed drawing
    \pgfmathsetmacro{\TotalW}{\cutX + \gapX + \rightW}
    \pgfmathsetmacro{\TotalH}{\cutY + \gapY + \bottomH}

    % Margins for outside labels (cells). Increase if your labels are longer.
    \newcommand{\padL}{2}   % left margin
    \newcommand{\padR}{2}   % right margin
    \newcommand{\padT}{2}   % top margin
    \newcommand{\padB}{2}   % bottom margin

    % Letter H coordinates
    \newcommand{\rowH}{4}
    \newcommand{\colH}{9}

    \begin{tikzpicture}[x=\cellw, y=-\cellh, line cap=rect] % negative y so (0,0) is top-left
      % Clamp bbox to include outside labels as well
      \path[use as bounding box]
        (-\padL, -\padT) rectangle (\TotalW+\padR, \TotalH+\padB);

      % Keep nodes from adding unexpected extra space
      \tikzset{every node/.style={outer sep=0pt}}

      % Cell coordinate
      \newcommand{\coord}[3]{%
        \node[font=\ttfamily\fontsize{2mm}{2mm}\selectfont] at (#2 + 0.5, #1 + 0.5) {#3};}

      % Monospace cell text centered in a cell
      \newcommand{\cell}[3]{%
        \node[font=\ttfamily\fontsize{4.4mm}{4.4mm}\selectfont,text=raspberry] at (#2 + 0.5, #1 + 0.5) {#3};}

      % Draw grid over integer-inclusive rectangle
      \newcommand{\gridrect}[4]{
        \pgfmathtruncatemacro{\xa}{#1}\pgfmathtruncatemacro{\ya}{#2}
        \pgfmathtruncatemacro{\xb}{#3}\pgfmathtruncatemacro{\yb}{#4}
        \foreach \x in {\xa,...,\xb}{\draw[gray!45, line width=0.15pt] (\x,\ya) -- (\x,\yb);}
        \foreach \y in {\ya,...,\yb}{\draw[gray!45, line width=0.15pt] (\xa,\y) -- (\xb,\y);}
      }

      % Curvy cutline
      \tikzset{cutline/.style={decorate, decoration={snake, amplitude=0.35mm, segment length=3.8mm}, draw=gray!45, thin}}

      % === HIGHLIGHT LAYER (drawn before the grid) ===
      % highlight H
      \fill[gray!10] (0,\rowH) rectangle (\colH+1,\rowH+1);
      \fill[gray!10] (\colH,0) rectangle (\colH+1,\rowH+1);

      % ===== TOP-LEFT (0..cutX, 0..cutY) =====
       \begin{scope}
          \clip
            (0,0) -- (\cutX,0)
            % Right (vertical) wavy edge: zero at y=0 and y=cutY
            -- plot[domain=0:\cutY, samples=\samps, variable=\t]
                 ({\cutX + \amp*sin(360*(\t - \cutY + 0.04)/\lam)}, {\t})
            % Bottom (horizontal) wavy edge: zero at x=cutX and x=0
            -- (\cutX,\cutY) % ensure we are exactly at the corner
            -- plot[domain=\cutX:0, samples=\samps, variable=\u]
                 ({\u}, {\cutY + \amp*sin(360*(\u - \cutX + 0.58)/\lam)})
            -- (0,\cutY) -- cycle;

          \gridrect{0}{0}{\cutX}{\cutY}
        \end{scope}

      % wavy edges
      \draw[cutline] (\cutX,\cutY) -- (\cutX,0) ;  % right wavy
      \draw[cutline] (\cutX,\cutY) -- (0,\cutY);  % bottom wavy

      % true outer borders
      \draw[thick] (0,0) -- (\cutX,0);    % top
      \draw[thick] (0,0) -- (0,\cutY);    % left

      % column coordinates up to H
      \foreach \x in {0,...,\colH}{\coord{-1}{\x}{\x};}
      \coord{-1}{\colH+1}{\raisebox{-1.25ex}{$\cdots$}}

      % row coordinates up to H
      \foreach \y in {0,...,\rowH}{\coord{\y}{-1}{\y};}
      \coord{\rowH+1}{-1}{$\vdots$}

      % ===== TOP-RIGHT (near top-right corner), compressed rightward =====
      % Size: rightW × (cutY+1 rows high)
      \begin{scope}[shift={(\cutX+\gapX, 0)}]
        \begin{scope}
            \clip
                (0,0) -- (\rightW,0) -- (\rightW,\cutY)                 % outer top & right
                % bottom wavy edge: from right->left, zero at u=0 (seam)
                -- plot[domain=\rightW:0, samples=\samps, variable=\u]
                     ({\u}, {\cutY + \amp*sin(360*(\u - 0 - 0.58)/\lam)})
                -- (0,\cutY) % ensure we are exactly at the corner
                % left wavy edge: from bottom->top, same phase as TL’s vertical
                -- plot[domain=\cutY:0, samples=\samps, variable=\t]
                     ({0 + \amp*sin(360*(\t - \cutY + 0.04)/\lam)}, {\t})
                -- cycle;

          \gridrect{0}{0}{\rightW}{\cutY}
        \end{scope}

        % wavy edges
        \draw[cutline] (0,\cutY) -- (0,0);        % faces vertical gap
        \draw[cutline] (0,\cutY) -- (\rightW,\cutY); % faces horizontal gap

        % true outer borders (top and right outer edge)
        \draw[thick] (0,0) -- (\rightW,0);
        \draw[thick] (\rightW,0) -- (\rightW,\cutY);
      \end{scope}

      % ===== BOTTOM-LEFT (near bottom-left corner), compressed downward =====
      % Size: (cutX+1 cols wide) × bottomH
      \begin{scope}[shift={(0, \cutY+\gapY)}]
        \begin{scope}
          \clip
            (0,\bottomH) -- (0,0)
            % top wavy edge: left->right, same phase as TL’s horizontal
            -- plot[domain=0:\cutX, samples=\samps, variable=\u]
                 ({\u}, {0 + \amp*sin(360*(\u - \cutX + 0.58)/\lam)})
            -- (\cutX,0) % ensure we are exactly at the corner
            % right wavy edge: top->bottom, zero at t=0 (seam)
            -- plot[domain=0:\bottomH, samples=\samps, variable=\t]
                 ({\cutX + \amp*sin(360*(\t - 0 - 0.04)/\lam)}, {\t})
            -- (\cutX,\bottomH) -- cycle;                        % outer bottom

          \gridrect{0}{0}{\cutX}{\bottomH}
        \end{scope}

        % wavy edges
        \draw[cutline] (\cutX,0) -- (0,0);        % faces horizontal gap
        \draw[cutline] (\cutX,0) -- (\cutX,\bottomH); % faces vertical gap

        % true outer borders (left and bottom outer edge)
        \draw[thick] (0,0) -- (0,\bottomH);
        \draw[thick] (0,\bottomH) -- (\cutX,\bottomH);
      \end{scope}

      % ===== BOTTOM-RIGHT (near bottom-right corner), compressed both ways =====
      \begin{scope}[shift={(\cutX+\gapX, \cutY+\gapY)}]
        \begin{scope}
          \clip
            (\rightW,\bottomH) -- (\rightW,0)
            % top wavy edge: right->left, zero at u=0 (seam with TR)
            -- plot[domain=\rightW:0, samples=\samps, variable=\u]
                 ({\u}, {0 + \amp*sin(360*(\u - 0 - 0.58)/\lam)})
            -- (0,0) % ensure we are exactly at the corner
            % left wavy edge: bottom->top (go up to close), zero at t=0 (seam with BL)
            -- plot[domain=0:\bottomH, samples=\samps, variable=\t]
                 ({0 + \amp*sin(360*(\t - 0 - 0.04)/\lam)}, {\t})
            -- (0,\bottomH) -- cycle;                           % outer bottom

          \gridrect{0}{0}{\rightW}{\bottomH}
        \end{scope}

        % wavy edges
        \draw[cutline] (0,0) -- (0,\bottomH);     % faces vertical gap
        \draw[cutline] (0,0) -- (\rightW,0);      % faces horizontal gap

        % true outer borders (right and bottom outer edges)
        \draw[thick] (\rightW,0) -- (\rightW,\bottomH);
        \draw[thick] (0,\bottomH) -- (\rightW,\bottomH);

      \end{scope}

      % ===== "Hello" at (rowH, colH) =====
      \cell{\rowH}{\colH + 0}{H}
      \cell{\rowH}{\colH + 1}{E}
      \cell{\rowH}{\colH + 2}{L}
      \cell{\rowH}{\colH + 3}{L}
      \cell{\rowH}{\colH + 4}{O}

    % ==== Corner labels ====
      \node[anchor=south east] at (0, 0) {\scriptsize (0,0)};
      \node[anchor=south west] at (\TotalW, 0) {\scriptsize (0,79)};
      \node[anchor=north east] at (0, \TotalH) {\scriptsize (59,0)};
      \node[anchor=north west] at (\TotalW, \TotalH) {\scriptsize (59,79)};

      % Axis captions
      \node at (\TotalW/2, -\padT/2 - 0.2) {\scriptsize 80 columns};
      \node[rotate=90] at (-\padT/2, \TotalH/2) {\scriptsize 60 rows};

    \end{tikzpicture}
} % end diagram
\end{document}
